% !TEX root = ../main_hessenberg.tex
\begin{comment}
\section{Compression as evidence of content}\label{sec:compression}

The criteria of \S\ref{sec:methodology} are internal: a green build,
and an apex signature that reads as the paper sentence. Both are
judgments the authors pass on their own artifact. This section adds
an external one. Aksenov, Bodnia, Freedman, and
Mulligan~\cite{aksenov2026compressionneedmodelingmathematics} take
compression as the measurable difference between mathematics and
mere formal text, and calibrate this on Mathlib: unfolding genuine
mathematics expands the source by many orders of magnitude, which
formal text restating its premises does not. We state their measure,
compute it for our development, and report where the development
lands.

\begin{definition}[Compression metrics,
{\cite[\S3]{aksenov2026compressionneedmodelingmathematics}}]\label{def:abfm}
\leanmarginLean{Util.Compression}{Util/Compression.lean}
Let $G$ be the directed graph whose vertices are the declarations of
a Lean development plus a synthetic vertex $\mathsf{Sort}$, with an
edge $u \to v$ of weight $w$ when the elaborated term of $u$ refers
to $v$ exactly $w$ times; declarations of the Lean core library and
$\mathsf{Sort}$ are \emph{primitive}. The \emph{wrapped length} of
$u$ is the number of tokens in its Lean source. Its \emph{unwrapped
length} is $|u|_G = \sum_i w_i \, |v_i|_G$ over the edges
$u \to v_i$, with $|u|_G = 1$ for primitives. Its \emph{depth} is
the length of the longest path in $G$ from $u$ to a primitive.
\end{definition}

Wrapped length is the cost of stating $u$ with the whole library at
hand; unwrapped length is the cost of stating it from primitives
alone. The ratio between the two measures the compression that the
intervening definitions buy, and depth counts how many layers of
definitions that compression is spread across.

We compute the three quantities for the thirteen public targets of
\texttt{HessenbergDigraphs} with a program written in Lean itself,
invoked as \texttt{lake exe hessenberg-compression}. The program is a
transcription of Definition~\ref{def:abfm}: it loads the compiled
library at runtime, collects the weighted reference multiset of each
elaborated term by structural recursion over the term (the
\texttt{collectElems} traversal of Aksenov et al.\ \S3.1), and evaluates
$|u|_G$ and depth by memoised recursion over $G$. Working inside Lean
gives the program the elaborated terms exactly as the kernel checked
them, rather than a reconstruction from printed output. Three
implementation choices on points that their description leaves
open (the tokenizer behind wrapped length, the references of
declarations that have a type but no term, such as inductive types
and their recursors, and cycle handling) are documented in the
source. The results appear in Table~\ref{tab:compression}.

\begin{table}[ht]
\centering
\renewcommand{\arraystretch}{1.05}
\begin{tabular}{@{}lrrr@{}}
\toprule
\textbf{Declaration} & \textbf{Wrapped} & \textbf{Unwrapped} & \textbf{Depth} \\
\midrule
\texttt{classification}              &     93 & $2.39 \times 10^{44}$ & 138 \\
\texttt{completeness}                &     61 & $2.38 \times 10^{44}$ & 137 \\
\texttt{digraph\_model}              &    123 & $4.95 \times 10^{41}$ & 114 \\
\texttt{exists\_givensFactorization} &    368 & $2.81 \times 10^{47}$ & 147 \\
\texttt{rigid}                       &  1\,201 & $2.44 \times 10^{12}$ &  28 \\
\texttt{singleton\_iso}              &    139 & $2.05 \times 10^{11}$ &  24 \\
\texttt{singleton\_iso\_unique}      &  1\,052 & $1.83 \times 10^{12}$ &  28 \\
\texttt{loopless\_iff}               &    566 & $1.01 \times 10^{8}$  &  21 \\
\texttt{count\_formula}              &     67 & $4.30 \times 10^{9}$  &  21 \\
\texttt{count\_loopless\_formula}    &     97 & $193\,677$            &   4 \\
\texttt{Arc}                         &     75 & $8$                   &   2 \\
\texttt{activeSet}                   &    100 & $1.01 \times 10^{37}$ & 101 \\
\texttt{givensProduct}               &     51 & $1.35 \times 10^{37}$ & 102 \\
\bottomrule
\end{tabular}
\caption{Compression metrics (wrapped, unwrapped, depth) of
Definition~\ref{def:abfm} for the thirteen public targets of
\texttt{HessenbergDigraphs}, computed by
\texttt{lake exe hessenberg-compression}.}
\label{tab:compression}
\end{table}

The four apex-level theorems and the foundation definition
\texttt{activeSet} have unwrapped lengths between $10^{37}$ and
$10^{47}$ and depths above $100$. This crosses the threshold by which
Aksenov et al.\ separate genuine hierarchical compression from mere
syntactic restatement, placing the project within the substantive-content
region of Mathlib: below their reported maximum of $10^{104}$, but
unambiguously inside the compression regime. The two enumeration
formulae and the singleton isomorphism sit in the
$10^{8}$--$10^{12}$ range at depths $20$--$28$, corresponding to
proofs that form a structural-lemma chain built on top of arithmetic
identities. \texttt{count\_loopless\_formula} is the smallest
target, at depth $4$ and unwrapped length $193\,677$: the proof is a
single closed-form arithmetic identity, which the metric reads off
directly.

This answers the question that the choice of a marshmallow invites:
a deliberately elementary problem has not left the artifact empty.
The internal criterion of \S\ref{sec:methodology} says the code reads
as mathematics; the measurement here says the case study of
\S\ref{sec:example} measures as mathematics.
\end{comment}
